% Cell-occupancy heatmap: substrate x cell counts at goal layer.
% Pure TikZ matrix (no pgfplots) so the annotation logic is robust.
% Designed for full text-width.
\begin{figure}[t]
\centering
\begin{tikzpicture}[
  font=\footnotesize,
  cell/.style={
    rectangle, draw=black!30, line width=0.3pt,
    minimum width=0.8cm, minimum height=0.6cm, align=center,
    inner sep=0pt
  },
]

% Cell list (16 columns)
\def\cells{1A,1B,2A,2B,3A,3B,4A,4B,5A,5B,6A,6B,7A,7B,8,9}
% Substrates (5 rows, top to bottom)
\def\substrates{FLT,SLT,FP,AS-v2,AS-v1}

% Counts data: row major.
% Order: FLT, SLT, FP, AS-v2, AS-v1.
% Cells: 1A 1B 2A 2B 3A 3B 4A 4B 5A 5B 6A 6B 7A 7B 8 9
% (zero entries shown as empty cells)
\def\dataFLT{0,0,2,0,0,0,3,0,0,0,0,1,2,0,7,93}
\def\dataSLT{0,0,0,0,0,0,0,0,0,0,0,0,0,0,17,78}
\def\dataFP{0,0,0,0,0,0,0,1,0,0,0,0,0,0,0,8}
\def\dataASvtwo{0,0,0,0,0,0,2,0,0,0,0,0,0,1,4,7}
\def\dataASvone{0,0,4,0,4,0,3,1,1,0,0,0,0,3,6,34}

% Header row: cell ids
\foreach \c [count=\i from 0] in 1A,1B,2A,2B,3A,3B,4A,4B,5A,5B,6A,6B,7A,7B,8,9 {
  \node[font=\scriptsize, black!70] at (\i*0.8 + 0.4, 0.7) {\c};
}

% Substrate row labels
\foreach \s [count=\j from 0] in FLT, SLT, FP, AS-v2, AS-v1 {
  \node[font=\scriptsize\bfseries, anchor=east] at (-0.1, -\j*0.6) {\s};
}

% FLT row
\foreach \v [count=\i from 0] in 0,0,2,0,0,0,3,0,0,0,0,1,2,0,7,93 {
  \pgfmathtruncatemacro{\shade}{min(75, \v)}
  \node[cell, fill=black!\shade] at (\i*0.8 + 0.4, 0) {%
    \ifnum\v>0\textcolor{black!90}{\v}\fi};
}
% SLT row
\foreach \v [count=\i from 0] in 0,0,0,0,0,0,0,0,0,0,0,0,0,0,17,78 {
  \pgfmathtruncatemacro{\shade}{min(75, \v)}
  \node[cell, fill=black!\shade] at (\i*0.8 + 0.4, -0.6) {%
    \ifnum\v>0\textcolor{black!90}{\v}\fi};
}
% FP row
\foreach \v [count=\i from 0] in 0,0,0,0,0,0,0,1,0,0,0,0,0,0,0,8 {
  \pgfmathtruncatemacro{\shade}{min(75, \v*8)}
  \node[cell, fill=black!\shade] at (\i*0.8 + 0.4, -1.2) {%
    \ifnum\v>0\textcolor{black!90}{\v}\fi};
}
% AS-v2 row
\foreach \v [count=\i from 0] in 0,0,0,0,0,0,2,0,0,0,0,0,0,1,4,7 {
  \pgfmathtruncatemacro{\shade}{min(75, \v*8)}
  \node[cell, fill=black!\shade] at (\i*0.8 + 0.4, -1.8) {%
    \ifnum\v>0\textcolor{black!90}{\v}\fi};
}
% AS-v1 row
\foreach \v [count=\i from 0] in 0,0,4,0,4,0,3,1,1,0,0,0,0,3,6,34 {
  \pgfmathtruncatemacro{\shade}{min(75, \v*2)}
  \node[cell, fill=black!\shade] at (\i*0.8 + 0.4, -2.4) {%
    \ifnum\v>0\textcolor{black!90}{\v}\fi};
}

% Axis labels
\node[font=\small\bfseries] at (6.4, 1.4) {Goal-layer cell};

\end{tikzpicture}
\caption{Goal-layer cell-occupancy across the five substrates (282 records,
counts shown). FLT and SLT concentrate in cell~9 (\texttt{fundamental\_result})
with a small tail in cell~8 (\texttt{supporting\_argument}). The AS-v1 and
AS-v2 substrates spread across cells 2A, 3A, 4A, 7B, 8, and 9 --- the
failure-mode cells the rubric exists to resolve. The First-Proof substrate
(FP) sits intermediate, with strong cell-9 occupancy and one cell-4B
(\texttt{boundary\_result}) entry on Problem~5. Empty cells are zero;
darker cells carry larger counts.}
\label{fig:cell_heatmap}
\end{figure}
